\newpage
\pagestyle{plain}

\section{ВЫВОДЫ}

В ходе выполнения работы получены следующие результаты:

\begin{enumerate}
	\item использован альтернативный способ разрешения системы условий порядка: в отличие от применяемых ранее подходов, где к системе применялся пакет символьных вычислений Maple, что приводило к большим временным затратам и громоздким формулам, в данной работе уравнения были решены численно с помощью библиотеки \texttt{scipy.optimize}.
	\item алгоритм стохастической оптимизации методом роя частиц был дополнен локальной полировкой симплекс-методом Нелдера--Мида, что привело к повышению точности решений.
	\item экспериментально зафиксировано кратное повышение точности разработанного метода по сравнению с классической непараметрической схемой Рунге--Кутты. Для линейного осциллятора погрешность снижена в $6842$ раза, для нелинейного уравнения Шрёдингера — в $424,38$ раза, для задачи двух тел — в $7.3$ раза при выраженной эллиптической орбите.
	\item эффективность изучаемого подхода подтверждена на примере задачи пяти внешних тел Солнечной системы, не имеющей аналитического решения. Алгоритм обеспечил долгосрочную стабильность моделей орбит планет.
\end{enumerate}